%%%%%%%%%%%%%%%%%
% CV tailored for Exakis Nelite - Senior Change Management Consultant
%%%%%%%%%%%%%%%%%

\documentclass[8pt,a4paper,ragged2e,withhyper]{../_shared/altacv}

\geometry{left=1.25cm,right=1.25cm,top=.5cm,bottom=1.cm,columnsep=1.2cm}

\usepackage{paracol}
\usepackage{fontawesome5}
\usepackage{hyperref}

\iftutex 
  \setmainfont{Roboto Slab}
  \setsansfont{Lato}
  \renewcommand{\familydefault}{\sfdefault}
\else
  \usepackage[rm]{roboto}
  \usepackage[defaultsans]{lato}
  \renewcommand{\familydefault}{\sfdefault}
\fi

\definecolor{SlateGrey}{HTML}{2E2E2E}
\definecolor{LightGrey}{HTML}{666666}
\definecolor{DarkPastelRed}{HTML}{8AC0D9}
\definecolor{PastelRed}{HTML}{00B0DA}
\definecolor{GoldenEarth}{HTML}{B4AD0C}

\colorlet{name}{black}
\colorlet{tagline}{PastelRed}
\colorlet{heading}{DarkPastelRed}
\colorlet{headingrule}{GoldenEarth}
\colorlet{subheading}{PastelRed}
\colorlet{accent}{PastelRed}
\colorlet{emphasis}{SlateGrey}
\colorlet{body}{LightGrey}

\definecolor{violet}{HTML}{5A2582}
\definecolor{rouge}{HTML}{F53B3B}
\definecolor{noir}{HTML}{00232C}
\definecolor{bleu}{HTML}{00B0DA}
\definecolor{rose}{HTML}{F831A0}
\definecolor{jaune}{HTML}{FFEC00}
\definecolor{vert}{HTML}{34AD6C}

\renewcommand{\namefont}{\Huge\rmfamily\bfseries}
\renewcommand{\personalinfofont}{\footnotesize}
\renewcommand{\cvsectionfont}{\LARGE\rmfamily\bfseries}
\renewcommand{\cvsubsectionfont}{\large\bfseries}
\renewcommand{\cvItemMarker}{{\small\textbullet}}
\renewcommand{\cvRatingMarker}{\faCircle}

\input{../_shared/pubs-num.tex}
\addbibresource{../_shared/sample.bib}


\begin{document}

\name{BOILEAU Guillaume}
\tagline{Senior Change Management Consultant | Digital Transformation, User Adoption, Training, Communication \& Governance}

\photoL{2.8cm}{../_shared/moi}

\personalinfo{%
  \email{guillaume.boileaupro@gmail.com}
  \phone{+33 6 59 94 85 84}
  \location{Alpes-Maritimes, FRANCE}
  \linkedin{guillaume-boileau-phd-891b81265}
  \researchgate{Guillaume-Boileau-2}
  \orcid{0000-0002-3576-6968}
}

\makecvheader
\vspace{-0.3cm}

\AtBeginEnvironment{itemize}{\small}
\columnratio{0.64}

\begin{paracol}{2}

\cvsection{Experience}

\cvevent{\textbf{Software Engineer - Digital Delivery, User Support \& Change Enablement}}{VEV Platform Services France}{2025 -- 2026}{Valbonne, France}

\textbf{Context:} Development and integration of software services for an EV charging CPMS platform connected to operational infrastructure, field equipment and business-facing services.

\textbf{Objective:} Support digital service delivery, operational adoption, incident follow-up and collaboration between technical stakeholders and users.

\begin{itemize}
  \item \textbf{Digital delivery:} Developed and maintained web and backend services using TypeScript, Angular and Node.js in an operational industrial environment.
  \item \textbf{Usage and process analysis:} Analysed existing practices around incidents, data exchanges, operational constraints, user feedback, task follow-up and technical deliverables.
  \item \textbf{Change support:} Structured corrective actions and user-oriented explanations from feedback, production issues, field-equipment constraints and data-flow anomalies.
  \item \textbf{Incident / issue management:} Investigated service behaviour, communication issues, interface interruptions and operational deviations affecting service continuity.
  \item \textbf{Data-flow reliability:} Implemented, maintained and monitored FTP-based operational data exchange services supporting industrial workflows.
  \item \textbf{Validation / acceptance:} Checked service behaviour, interface continuity and data consistency in production-like or operational conditions.
  \item \textbf{Collaborative tooling:} Used Zenhub to support backlog follow-up, action tracking, iterative delivery and technical coordination.
  \item \textbf{Documentation and support:} Formalised technical findings, open points, corrective actions and usage-oriented explanations to improve shared understanding and follow-up.
  \item \textbf{Stakeholder coordination:} Worked with technical contributors to clarify priorities, unblock issues and support pragmatic delivery decisions.
\end{itemize}

\divider

\cvevent{\textbf{Senior R\&D Engineer - Digital Transformation, Methods \& Stakeholder Enablement}}{CNES, Laboratoire Lagrange, Observatoire de la Côte d'Azur}{2023 -- 2025}{Nice, France}

\textbf{Context:} Engineering activities for the LISA ESA/NASA space mission, involving complex software chains, model integration, simulation workflows, validation campaigns and international coordination.

\textbf{Objective:} Structure, document and improve collaborative methods for complex data/software workflows while supporting adoption, stakeholder alignment and evidence-based decision-making.

\begin{itemize}
  \item \textbf{Transformation and methods:} Designed and improved working methods for model integration, comparison campaigns, validation workflows and evidence consolidation.
  \item \textbf{Functional coordination:} Coordinated technical activities involving contributors, research engineers, technicians and Master's thesis interns, with progress follow-up, stakeholder alignment and deliverable consolidation.
  \item \textbf{Governance logic:} Defined validation logic, expected outputs, review practices, acceptance criteria and reproducible workflow principles.
  \item \textbf{Tooling and workflows:} Developed CATpyLISA, a Python platform for large-scale model integration, comparison, validation and technical review.
  \textcolor{DarkPastelRed}{\href{https://gitlab.oca.eu/gboileau/lisa_l3_tools}{\bf GitLab: CATpyLISA}}
  \item \textbf{Documentation:} Used Confluence to support technical documentation, coordination notes, validation follow-up, decision logs and knowledge sharing.
  \item \textbf{KPI / indicators logic:} Consolidated validation status, comparison results, deviations, open points and progress material to support steering and technical decisions.
  \item \textbf{Continuous improvement:} Identified inconsistencies, bottlenecks, quality issues and technical risks affecting reliability, progress or delivery quality.
  \item \textbf{Reporting:} Produced technical notes, validation summaries, analysis reports, progress material and review supports for contributors and stakeholders.
  \item \textbf{Change enablement:} Helped stakeholders adopt shared workflows, common documentation practices, clearer validation / review methods and reusable support material.
  \item \textbf{International environment:} Operated in a collaborative framework requiring technical English, structured communication and versioned workflows.
\end{itemize}

\divider

\cvevent{\textbf{Data Scientist - Data Adoption, Technical Analysis \& Reporting}}{Universiteit Antwerpen, Belgium}{2021 -- 2023}{Antwerp, Belgium}

\textbf{Context:} Data analysis and performance studies for precision instruments in a high-requirement international environment linked to gravitational-wave detector performance.

\textbf{Objective:} Develop reproducible analysis workflows, support data-driven decision-making and communicate complex technical conclusions to international stakeholders.

\begin{itemize}
  \item \textbf{Data workflow design:} Developed Python-based pipelines for noise source identification, uncertainty estimation and statistical modelling.
  \item \textbf{Process reliability:} Improved reproducibility, traceability and robustness of technical analyses through structured methods and documented workflows.
  \item \textbf{Audit-like analysis:} Analysed instrumental and environmental effects impacting system sensitivity, stability and measurement quality.
  \item \textbf{KPI-style reporting:} Synthesised performance indicators, deviations and dominant contributors into actionable technical conclusions.
  \item \textbf{Documentation:} Produced reproducible and traceable technical analyses for international engineering and research stakeholders.
  \item \textbf{Communication and pedagogy:} Communicated complex technical results through structured reporting, presentations and collaborative review.
\end{itemize}

\divider

\cvevent{\textbf{Junior R\&D Engineer - Simulation, Software Workflows \& Validation}}{ARTEMIS, Observatoire de la Côte d'Azur}{2018 -- 2021}{Nice, France}

\textbf{Context:} Engineering and analysis activities for gravitational-wave mission preparation, involving signal analysis, simulation, software workflows and noise characterization.

\textbf{Objective:} Build robust analysis methods and validation practices for complex signals and demanding system environments.

\begin{itemize}
  \item \textbf{Workflow formalisation:} Built simulation approaches for complex signal environments, uncertainty studies and large-scale datasets.
  \item \textbf{Validation methods:} Compared simulated outputs, model behaviours and analysis results to assess consistency, robustness and reliability.
  \item \textbf{Technical investigations:} Investigated noise sources through cross-sensor correlation analyses in diagnostic contexts.
  \item \textbf{Continuous improvement mindset:} Identified limiting factors affecting mission-level performance and proposed analysis improvements.
  \item \textbf{Scientific computing:} Developed strong expertise in Python-based analysis, modelling, documentation and reproducible workflows.
\end{itemize}

\switchcolumn


\cvsection{Profile for Exakis Nelite}

\textbf{Digital transformation and change-management profile with strong experience in structuring methods, analysing user and stakeholder needs, documenting processes, producing support material and helping multidisciplinary teams adopt clearer ways of working}. Background developed through international scientific, software and industrial environments, reinforced by operational digital delivery experience.

For a \textbf{Senior Change Management Consultant} role, I bring a hybrid profile combining technical depth, communication, pedagogy and transformation support: framing needs, analysing impacts, mapping use cases, structuring governance and adoption trajectories, producing documentation and training-style material, facilitating workshops and supporting users or stakeholders through complex digital changes.

\begin{itemize}
  \item Strong fit with change-management activities: needs analysis, impact assessment, use-case mapping, adoption support, documentation and communication material.
  \item Experience translating complex technical topics into clear explanations, support material, decision aids and stakeholder-facing documentation.
  \item Ability to interact with technical experts, operational contributors, project stakeholders and decision-makers in international contexts.
  \item Practical experience supporting new workflows, shared practices, governance logic and continuous improvement in software/data environments.
  \item Comfortable with digital transformation topics around data, AI, software tools, collaborative workflows, user support and operational adoption.
  \item Strong written and oral communication, with a rigorous and pedagogical approach to complex heterogeneous environments.
\end{itemize}

\cvsection{Fit for Change Management}

\begin{itemize}
  \item Relevant background for framing transformation needs, analysing impacts, structuring target practices and defining pragmatic adoption trajectories.
  \item Strong experience producing documentation, guides, review material, communication supports, reporting summaries and training-style explanations.
  \item Experience facilitating alignment across stakeholders through workshops, technical reviews, progress meetings and collaborative decision-making.
  \item Practical understanding of digital environments: applications, APIs, data flows, Python, SQL, Git/GitLab, Jira-style tools, Confluence, Linux, CI/CD and monitoring.
  \item Strong fit with governance and usage strategies through process definition, documentation standards, KPI follow-up and continuous-improvement loops.
  \item Able to bridge business needs, user constraints and technical implementation topics in data, AI, software and operational IT contexts.
\end{itemize}

\cvsection{Microsoft \& Digital Awareness}

\begin{itemize}
  \item Familiar with professional digital collaboration environments, Office tools, documentation practices, reporting supports and user-oriented communication.
  \item Strong ability to ramp up on Microsoft ecosystem topics such as Microsoft 365, Azure, Power Platform, Copilot / AI usage and cloud-oriented transformation contexts.
  \item Relevant technical culture in AI, data, software engineering, cybersecurity awareness and cloud/infrastructure concepts.
\end{itemize}

\cvsection{Languages}

\cvskill{English}{4}
\divider
\cvskill{Spanish}{2}
\divider

\medskip

\cvsection{Education}

\cvevent{PhD in Physics}{Université Côte d'Azur}{2018 -- 2021}{Nice, France}

\divider

\cvevent{Selective Graduate Program in Fundamental Physics (Magistère)}{Université Paris-Saclay}{2016 -- 2018}{Orsay, France}

\divider

\cvevent{MSc in Astronomy and Astrophysics}{Observatoire de Paris / Université Paris-Saclay}{2017 -- 2018}{Paris, France}

\divider

\cvevent{BSc in Fundamental Physics}{Université Paris-Saclay}{2015 -- 2016}{Orsay, France}

\end{paracol}

\cvsection{Professional Training}

\cvevent{Machine Learning in Python with scikit-learn}{FUN MOOC -- Inria}{2026}{France Université Numérique}

\small{Predictive modelling, supervised learning, model evaluation, cross-validation, metrics, pipelines and practical machine learning workflows with Python and scikit-learn.}

\divider

\cvevent{Recherche reproductible : principes méthodologiques pour une science transparente}{FUN MOOC -- Inria}{2025}{France Université Numérique}

\small{Reproducible research, GitLab, notebooks, data management, computational documentation, reproducible software environments and methodological transparency.}

\divider

\cvevent{Supervised Machine Learning (Regression \& Classification)}{DeepLearning.AI - Andrew Ng}{2020}{Coursera}

\divider

\cvevent{Continuous Professional Training -- Software, Data, AI and Systems}{OpenClassrooms}{2024 -- 2026}{}

\small{Python, SQL, Machine Learning, AI, Linux, JavaScript, TypeScript, Angular, Node.js, Django, REST APIs, data analysis, algorithms, C/C++ and software engineering fundamentals.}

\cvsection{Projects}

\cvevent{CNES / LISA -- Digital Methods, Change Enablement \& Validation Workflows}{ESA/NASA Space Mission - Large-scale International Collaboration}{2023 -- 2025}{Nice, France}

Structuring of software, data-processing and validation workflows for the LISA space mission within a large international engineering and scientific collaboration, with strong focus on adoption, documentation and stakeholder alignment.

\textbf{Change-management relevance:} Analysis of existing practices, workflow formalisation, contributor coordination, adoption support, documentation, action follow-up, impact analysis and stakeholder alignment.

\textbf{Digital transformation relevance:} Work involving complex software workflows, GitLab-based development, reproducible environments, technical documentation, validation evidence, indicators and versioned deliverables.

\textbf{Scale:} Major international space mission with an estimated budget of approximately 2 billion EUR over 10 years.

\divider

\cvevent{Einstein Telescope / LIGO--Virgo--KAGRA -- Technical Studies \& Collaborative Project Workflows}{International Collaborations}{2021 -- 2023}{Antwerp, Belgium}

Contribution to collaborative developments in data analysis, software methods and technical investigations within large-scale international scientific and engineering collaborations.

\textbf{Transformation relevance:} Documentation of technical studies, progress reporting, collaborative review, prioritisation of analysis tasks and communication with international stakeholders.

\textbf{System impact:} Studies on environmental and instrumental noise sources supporting the identification of performance limitations and design constraints for next-generation gravitational-wave observatories.


\cvsection{Strengths}

\vspace{-.3cm}

\cvsubsection{Change management, adoption \& transformation}

{\small
\cvtag{Change Management}
\cvtag{Digital Transformation}
\cvtag{User Adoption}
\cvtag{Impact Analysis}
\cvtag{Use-Case Mapping}
\cvtag{Target Operating Model Awareness}
\cvtag{Adoption Roadmaps}
\cvtag{Communication Plans}
\cvtag{Training Plans}
\cvtag{Governance}
\cvtag{Continuous Improvement}
\cvtag{KPI Follow-Up}
\cvtag{Stakeholder Alignment}
\cvtag{Knowledge Sharing}
\cvtag{Process Documentation}
\cvtag{Action Plans}
}

\divider\smallskip
\vspace{-.3cm}


\cvsubsection{Workshops, governance \& project tooling}

{\small
\cvtag{Agile}
\cvtag{Scrum Awareness}
\cvtag{Kanban Awareness}
\cvtag{SAFe Awareness}
\cvtag{Workshop Facilitation}
\cvtag{User Workshops}
\cvtag{Backlog Follow-Up}
\cvtag{JIRA}
\cvtag{Confluence}
\cvtag{Zenhub}
\cvtag{Trello Awareness}
\cvtag{Git / GitLab}
\cvtag{GitHub}
\cvtag{Project Coordination}
\cvtag{Stakeholder Alignment}
\cvtag{Cross-Functional Coordination}
\cvtag{Progress Reporting}
}

\divider\smallskip
\vspace{-.3cm}

\cvsubsection{Digital, data, AI \& technical environment}

{\small
\cvtag{Information Systems}
\cvtag{Applications}
\cvtag{Software Workflows}
\cvtag{Data Flows}
\cvtag{APIs}
\cvtag{REST API}
\cvtag{Python}
\cvtag{SQL}
\cvtag{Linux}
\cvtag{Docker}
\cvtag{CI/CD}
\cvtag{TypeScript}
\cvtag{Angular}
\cvtag{Node.js}
\cvtag{Operational Monitoring}
\cvtag{OpenSearch}
\cvtag{Power BI Awareness}
\cvtag{Microsoft 365 Awareness}
\cvtag{Azure Awareness}
\cvtag{AI Usage Awareness}
}

\divider\smallskip
\vspace{-.3cm}

\cvsubsection{Documentation, communication \& training}

{\small
\cvtag{Technical Documentation}
\cvtag{User Guides}
\cvtag{Training Material}
\cvtag{Communication Supports}
\cvtag{Adoption Supports}
\cvtag{Workshop Material}
\cvtag{Pedagogy}
\cvtag{Executive Presentations}
\cvtag{Decision Logs}
\cvtag{Meeting Notes}
\cvtag{Synthesis}
\cvtag{Technical English}
\cvtag{French Professional Writing}
\cvtag{Stakeholder Communication}
}

\divider\smallskip
\vspace{-.3cm}

\cvsubsection{Analysis, communication \& soft skills}

{\small
\cvtag{Analytical Thinking}
\cvtag{Problem Solving}
\cvtag{Root-Cause Analysis}
\cvtag{Rigor}
\cvtag{Autonomy}
\cvtag{Adaptability}
\cvtag{Team Spirit}
\cvtag{Clear Communication}
\cvtag{PowerPoint}
\cvtag{Excel}
\cvtag{Word}
\cvtag{MS Office}
\cvtag{Executive Presentations}
\cvtag{Reporting Dashboards}
}

\end{document}
